
For the fixed birational map $L(u,v)=(u,(u^2+1)v)$, reduction at every
prime $p\equiv3\pmod4$ permutes all of $\Fp^2$. Its point-weighted
distribution of periods divided by $p$ has no limit. At every fixed
$1/2\le x<1$, the distribution function is
$1-\varphi(p-1)/(p-1)+o(1)$, with limsup $1$ and liminf at most $23/24$.
Moreover, neither these distribution-function values nor the period
measures converge on any further prime subset of positive upper relative
density. The latter statement uses the classical continuous empirical
law for $\varphi(p-1)/(p-1)$. A supplementary proposition treats all odd
primes on the invariant regular domain, and under uniform-point deletion
or arbitrary exceptional-period completion. These conventions cannot
restore convergence on any positive-upper-density prime selection.
The contribution is a precisely delimited
application to the broad existence assertion in the 2004 Oberwolfach
problem on finite-field maps. The triangular family, primitive-value
estimates, totient mean and scalar distribution law are established
ingredients; no new arithmetic theorem or historical firstness is claimed.
